\chapter{Esercizi per argomento}
\label[appendix]{chap:esercizi-argomenti}

\section{Limiti}
\label{sec:esercizi-limiti}

\subsection{Semplificazioni algebriche}
\label{sec:esercizi-limiti-semplificazioni-algebriche}

\subsection{Limiti unilateri e vincoli di dominio}
\label{sec:esercizi-limiti-limiti-unilateri-e-dominio}

\subsection{Riconduzione ai limiti notevoli}
\label{sec:esercizi-limiti-limiti-notevoli}

\subsection{Cambio di variabile nei limiti}
\label{sec:esercizi-limiti-cambio-di-variabile}

\subsection{Razionalizzazione}
\label{sec:esercizi-limiti-razionalizzazione}

\subsection{Trasformazioni mediante identità trigonometriche}
\label{sec:esercizi-limiti-identita-trigonometriche}

\subsection{Equivalenze asintotiche}
\label{sec:esercizi-limiti-equivalenze-asintotiche}

\subsection{Sviluppi di Taylor e cancellazioni nei limiti}
\label{sec:esercizi-limiti-sviluppi-di-taylor-e-cancellazioni}

\subsection{Forme esponenziali e passaggio al logaritmo}
\label{sec:esercizi-limiti-forme-esponenziali-e-logaritmi}

\subsection{Regola di de l'Hôpital}
\label{sec:esercizi-limiti-de-lhopital}

\subsection{Teorema del confronto e stime}
\label{sec:esercizi-limiti-teorema-del-confronto-e-stime}

\subsection{Confronto dei termini nei limiti all'infinito}
\label{sec:esercizi-limiti-limiti-all-infinito}

\subsection{Passaggio dai limiti di funzioni ai limiti di successioni}
\label{sec:esercizi-limiti-limiti-di-successioni}

\subsection{Limiti di espressioni contenenti integrali}
\label{sec:esercizi-limiti-limiti-con-integrali}

\subsection{Discussione dei limiti con parametri}
\label{sec:esercizi-limiti-limiti-con-parametri}


\clearpage
\section{Studio di funzione}
\label{sec:esercizi-studio-di-funzione}

\subsection{Determinazione del dominio}
\label{sec:esercizi-studio-di-funzione-dominio}

\subsection{Simmetrie e periodicità}
\label{sec:esercizi-studio-di-funzione-simmetrie-e-periodicita}

\subsection{Studio del segno mediante gli zeri}
\label{sec:esercizi-studio-di-funzione-segno-e-zeri}

\subsection{Localizzazione degli zeri e delle intersezioni}
\label{sec:esercizi-studio-di-funzione-zeri-e-intersezioni}

\subsection{Limiti agli estremi del dominio}
\label{sec:esercizi-studio-di-funzione-limiti-agli-estremi-del-dominio}

\subsection{Ricerca degli asintoti}
\label{sec:esercizi-studio-di-funzione-asintoti}

\subsection{Continuità e prolungamenti continui}
\label{sec:esercizi-studio-di-funzione-continuita-e-prolungamenti}

\subsection{Valore assoluto e funzioni a tratti}
\label{sec:esercizi-studio-di-funzione-valore-assoluto-e-funzioni-a-tratti}

\subsection{Applicazione delle regole di derivazione}
\label{sec:esercizi-studio-di-funzione-regole-di-derivazione}

\subsection{Limiti della derivata nei punti singolari}
\label{sec:esercizi-studio-di-funzione-limiti-della-derivata}

\subsection{Classificazione dei punti di non derivabilità}
\label{sec:esercizi-studio-di-funzione-punti-di-non-derivabilita}

\subsection{Studio della monotonia mediante la derivata}
\label{sec:esercizi-studio-di-funzione-monotonia}

\subsection{Punti critici, estremi locali ed estremi assoluti}
\label{sec:esercizi-studio-di-funzione-punti-critici-ed-estremi}

\subsection{Estremo superiore e inferiore dell'immagine}
\label{sec:esercizi-studio-di-funzione-estremo-superiore-e-inferiore}

\subsection{Derivata seconda, convessità e flessi}
\label{sec:esercizi-studio-di-funzione-derivata-seconda-e-convessita}

\subsection{Costruzione del grafico qualitativo}
\label{sec:esercizi-studio-di-funzione-grafico-qualitativo}

\subsection{Studio di funzioni con parametri}
\label{sec:esercizi-studio-di-funzione-funzioni-con-parametri}


\clearpage
\section{Successioni}
\label{sec:esercizi-successioni}

\subsection{Limiti mediante le gerarchie di crescita}
\label{sec:esercizi-successioni-limiti-e-gerarchie}

\subsection{Determinazione dell'ordine asintotico}
\label{sec:esercizi-successioni-ordine-asintotico}

\subsection{Determinazione del segno definitivo}
\label{sec:esercizi-successioni-segno-definitivo}

\subsection{Fattoriali e rapporti di termini consecutivi}
\label{sec:esercizi-successioni-fattoriali-e-rapporti}

\subsection{Limiti di radici di indice variabile}
\label{sec:esercizi-successioni-radici-n-esime}

\subsection{Trasformazione delle differenze logaritmiche}
\label{sec:esercizi-successioni-differenze-logaritmiche}

\subsection{Sviluppi di Taylor nella variabile \(1/n\)}
\label{sec:esercizi-successioni-taylor-in-uno-su-n}

\subsection{Successioni alterne e sottosuccessioni}
\label{sec:esercizi-successioni-successioni-alterne}

\subsection{Monotonia e limitatezza delle successioni}
\label{sec:esercizi-successioni-monotonia-e-limitatezza}

\subsection{Studio delle successioni ricorsive}
\label{sec:esercizi-successioni-successioni-ricorsive}

\subsection{Dallo studio della successione al criterio per la serie}
\label{sec:esercizi-successioni-dalla-successione-alla-serie}


\clearpage
\section{Serie numeriche}
\label{sec:esercizi-serie}

\subsection{Verifica della condizione necessaria di convergenza}
\label{sec:esercizi-serie-condizione-necessaria}

\subsection{Riconduzione alle serie geometriche}
\label{sec:esercizi-serie-serie-geometriche}

\subsection{Confronto con serie armoniche generalizzate e logaritmiche}
\label{sec:esercizi-serie-serie-armoniche-generalizzate-e-logaritmi}

\subsection{Scomposizione telescopica}
\label{sec:esercizi-serie-serie-telescopiche}

\subsection{Criterio del confronto diretto}
\label{sec:esercizi-serie-criterio-del-confronto}

\subsection{Criterio del confronto asintotico}
\label{sec:esercizi-serie-confronto-asintotico}

\subsection{Criterio del rapporto}
\label{sec:esercizi-serie-criterio-del-rapporto}

\subsection{Criterio della radice}
\label{sec:esercizi-serie-criterio-della-radice}

\subsection{Serie con fattoriali e potenze}
\label{sec:esercizi-serie-fattoriali-e-potenze}

\subsection{Sviluppi di Taylor del termine generale}
\label{sec:esercizi-serie-taylor-nel-termine-generale}

\subsection{Stime dei termini e dei resti}
\label{sec:esercizi-serie-stime-e-maggiorazioni}

\subsection{Convergenza assoluta e convergenza semplice}
\label{sec:esercizi-serie-convergenza-assoluta-e-semplice}

\subsection{Criterio di Leibniz}
\label{sec:esercizi-serie-criterio-di-leibniz}

\subsection{Serie con fattori oscillanti}
\label{sec:esercizi-serie-serie-oscillanti}

\subsection{Discussione della convergenza con parametri}
\label{sec:esercizi-serie-serie-a-parametri}

\subsection{Parametri negli argomenti logaritmici e nelle basi}
\label{sec:esercizi-serie-parametri-in-logaritmi-e-basi}

\subsection{Raggio di convergenza e studio degli estremi}
\label{sec:esercizi-serie-serie-di-potenze}


\clearpage
\section{Integrali}
\label{sec:esercizi-integrali}

\subsection{Riconoscimento delle primitive immediate}
\label{sec:esercizi-integrali-primitive-immediate}

\subsection{Riconoscimento della derivata interna}
\label{sec:esercizi-integrali-riconoscimento-derivata-interna}

\subsection{Integrazione per sostituzione}
\label{sec:esercizi-integrali-integrazione-per-sostituzione}

\subsection{Integrazione per parti}
\label{sec:esercizi-integrali-integrazione-per-parti}

\subsection{Completamento del quadrato}
\label{sec:esercizi-integrali-completamento-del-quadrato}

\subsection{Funzioni razionali e decomposizione in fratti semplici}
\label{sec:esercizi-integrali-funzioni-razionali-e-fratti-semplici}

\subsection{Identità trigonometriche per l'integrazione}
\label{sec:esercizi-integrali-identita-trigonometriche}

\subsection{Sostituzione di Weierstrass}
\label{sec:esercizi-integrali-sostituzione-di-weierstrass}

\subsection{Sostituzioni algebriche e trigonometriche con radicali}
\label{sec:esercizi-integrali-sostituzioni-con-radicali}

\subsection{Sostituzioni iperboliche nei radicali}
\label{sec:esercizi-integrali-sostituzione-iperbolica}

\subsection{Valore assoluto e suddivisione dell'intervallo}
\label{sec:esercizi-integrali-valore-assoluto-e-suddivisione}

\subsection{Simmetrie e parità negli integrali}
\label{sec:esercizi-integrali-simmetrie-e-parita}

\subsection{Teorema fondamentale del calcolo e funzioni integrali}
\label{sec:esercizi-integrali-teorema-fondamentale-del-calcolo}

\subsection{Sviluppi di Taylor delle funzioni integrali}
\label{sec:esercizi-integrali-taylor-di-funzioni-integrali}

\subsection{Riconoscimento e separazione degli integrali impropri}
\label{sec:esercizi-integrali-riconoscere-integrali-impropri}

\subsection{Distinzione tra calcolo e convergenza}
\label{sec:esercizi-integrali-calcolo-e-convergenza}

\subsection{Confronto con integrali modello di potenze}
\label{sec:esercizi-integrali-integrali-modello-di-potenze}

\subsection{Confronto asintotico per integrali impropri}
\label{sec:esercizi-integrali-confronto-asintotico-impropri}

\subsection{Analisi delle singolarità logaritmiche}
\label{sec:esercizi-integrali-singolarita-logaritmiche}

\subsection{Confronti con esponenziali all'infinito}
\label{sec:esercizi-integrali-esponenziali-all-infinito}

\subsection{Convergenza degli integrali impropri con parametri}
\label{sec:esercizi-integrali-integrali-impropri-con-parametri}


\clearpage
\section{Equazioni differenziali}
\label{sec:esercizi-equazioni-differenziali}

\subsection{Classificazione dell'equazione e scelta del metodo}
\label{sec:esercizi-equazioni-differenziali-classificazione-equazioni}

\subsection{Ricerca e conservazione delle soluzioni di equilibrio}
\label{sec:esercizi-equazioni-differenziali-soluzioni-di-equilibrio}

\subsection{Equazioni a variabili separabili}
\label{sec:esercizi-equazioni-differenziali-equazioni-a-variabili-separabili}

\subsection{Equazioni separabili con decomposizione in fratti semplici}
\label{sec:esercizi-equazioni-differenziali-separabili-con-fratti-semplici}

\subsection{Equazioni lineari del primo ordine e fattore integrante}
\label{sec:esercizi-equazioni-differenziali-equazioni-lineari-primo-ordine}

\subsection{Equazioni lineari omogenee del secondo ordine a coefficienti costanti}
\label{sec:esercizi-equazioni-differenziali-equazioni-lineari-omogenee-secondo-ordine}

\subsection{Radice caratteristica doppia}
\label{sec:esercizi-equazioni-differenziali-radici-caratteristiche-multiple}

\subsection{Soluzione particolare con termine noto polinomiale}
\label{sec:esercizi-equazioni-differenziali-termine-noto-polinomiale}

\subsection{Soluzione particolare con termine noto esponenziale}
\label{sec:esercizi-equazioni-differenziali-termine-noto-esponenziale}

\subsection{Parametri e risonanza nelle equazioni lineari}
\label{sec:esercizi-equazioni-differenziali-parametri-e-risonanza}

\subsection{Imposizione dei dati nel problema di Cauchy}
\label{sec:esercizi-equazioni-differenziali-problema-di-cauchy}

\subsection{Possibile non unicità nelle equazioni con radicali}
\label{sec:esercizi-equazioni-differenziali-non-unicita-con-radicali}

\subsection{Determinazione dell'intervallo massimale}
\label{sec:esercizi-equazioni-differenziali-intervallo-massimale}

\subsection{Verifica diretta delle soluzioni}
\label{sec:esercizi-equazioni-differenziali-verifica-delle-soluzioni}


\clearpage
\section{Taylor e tecniche asintotiche}
\label{sec:esercizi-taylor-e-asintotici}

\subsection{Scelta tra equivalenze, stime e sviluppi di Taylor}
\label{sec:esercizi-taylor-e-asintotici-scelta-del-metodo-di-applicazione}

\subsection{Uso della notazione \(o\) piccolo e \(O\) grande}
\label{sec:esercizi-taylor-e-asintotici-notazione-o-piccolo-e-o-grande}

\subsection{Impiego degli sviluppi notevoli}
\label{sec:esercizi-taylor-e-asintotici-sviluppi-notevoli}

\subsection{Sviluppi centrati in un punto assegnato}
\label{sec:esercizi-taylor-e-asintotici-sviluppi-centrati-in-un-punto}

\subsection{Sviluppi binomiali per potenze reali}
\label{sec:esercizi-taylor-e-asintotici-potenze-reali-e-binomiali}

\subsection{Ordine di sviluppo dopo le prime cancellazioni}
\label{sec:esercizi-taylor-e-asintotici-ordine-della-prima-cancellazione}

\subsection{Composizione degli sviluppi}
\label{sec:esercizi-taylor-e-asintotici-composizione-di-sviluppi}

\subsection{Prodotti e quozienti di sviluppi}
\label{sec:esercizi-taylor-e-asintotici-prodotti-e-quozienti}

\subsection{Cancellazioni dei coefficienti con parametri}
\label{sec:esercizi-taylor-e-asintotici-cancellazioni-con-parametri}

\subsection{Segno definitivo dal primo termine non nullo}
\label{sec:esercizi-taylor-e-asintotici-segno-definitivo-dal-primo-termine}

\subsection{Sviluppi delle radici di indice variabile tramite logaritmi}
\label{sec:esercizi-taylor-e-asintotici-radici-n-esime-con-logaritmi}

\subsection{Confronto asintotico di potenze e logaritmi}
\label{sec:esercizi-taylor-e-asintotici-confronto-di-potenze-e-logaritmi}

